% SPDX-License-Identifier: Apache-2.0
\section{Pins, and an Interrupt on One}
\label{sec:x-gpio}

A pin is the least interesting thing a chip has and the first thing
anybody wants. Until there is one, a running program can say nothing
except through the serial port, and a board's four LEDs and its keys
are wire that goes nowhere.

\code{Gpio<N>} is \code{N} of them behind AXI-Lite, in seven words: what
to drive, what the pins read, which way each pin faces, and then the
four that make an interrupt, an enable, a kind, a polarity and a
status. The part drives no pin itself. It offers the three wires a pad
wants, the value, the direction and the value read back, and a board's
top level joins them to a tristate buffer. That is not a detour: a pad
in this language carries nothing until there is a netlist, so a part
written around one is a part no simulation can run, and a peripheral
nobody can simulate is a peripheral nobody can check.

An input arrives through two flip-flops before anything reads it.
A pin is a wire to the outside and its edges fall where they like, so
a value taken straight into logic can be half way between one and zero
when two gates read it, and the two gates can disagree. Two flip-flops
is the usual price of not caring.

The interrupt is per pin and is four registers rather than one,
because there are four questions and they are independent: whether the
pin may interrupt at all, whether it is a level or an edge, which
level or which edge, and whether it has happened. The status bit is
cleared by writing a one to it and not a zero, so two programs
clearing different pins cannot lose each other's work. A pin set for a
level sets its bit again in the cycle after it is cleared while the
level lasts, which is what a level means, and the example shows that
happening rather than asserting it.

The run is a bring-up program. It makes the low four pins outputs and
drives a pattern, reads the high four back as the testbench moves
them, asks for an interrupt on one pin's rising edge, and watches the
line rise on the rise and stay down on the fall. Then it clears the
status and watches the line fall, switches the pin to a level, clears
again, and finds the bit back.

One detail of the testbench is worth the sentence it takes. The pins
are moved between one cycle and the next, from the loop, and not from
a process inside the simulation. A process that set them during the
cycle would race the peripheral reading them, and the netlist, whose
port is settled before the edge, would then disagree with the run.
The check would fail, and it would be the testbench that was wrong.

\begin{figure}[htbp]
\centering
\input{gpio_timing.tex}
\caption{The pins and the interrupt: the value driven out, the value
read in two cycles behind the pin, the status bit set by the rising
edge and cleared by the write, and the line that follows it.}
\label{fig:x-gpio}
\end{figure}

\lstinputlisting[caption={\code{ex\_gpio.rs}. Run with \code{bazel run //lib/examples:ex\_gpio}.},label={lst:x-gpio}]{lib/examples/ex_gpio.rs}

\lstinputlisting[style=ascii,caption={What \code{ex\_gpio} printed when the build ran it: what the program read at each step, and then the peripheral's Verilog.},label={lst:x-gpio-out}]{out_gpio.txt}
